\clearpage\SourceLeaf{427}\NoSourceHead
\begingroup\makeatletter\let\veqno\@@leqno\makeatother
\newcommand{\vsroot}[1]{\mathord{\vcenter{\hbox{\Large\(\surd\)}}}#1}
\begin{center}[VII.]\par\medskip\large [ÜBER DIE KONVERGENZ\\DER ENTWICKLUNG DER MITTELPUNKTSGLEICHUNG.]\end{center}
\begin{center}\rule{.36\linewidth}{.4pt}\\[2pt]{\small [Vier Zettel, Astr. d{\scriptsize 1}, Kapsel {\scriptsize 90}.]}\\[-2pt]\rule{.36\linewidth}{.4pt}\end{center}
\begin{center}[1.]\end{center}
\noindent{\small [Erster Zettel]}\par
\noindent\hspace*{\parindent}[Es sei für]
\begin{equation}t=\cos u+i\sin u\tag*{[1]}\end{equation}
\begin{equation}\frac{1}{(1-\alpha t)(1-\alpha t^{-1})}=T,\tag*{[2]}\end{equation}
\noindent so haben [wir]
\begin{equation}(1-\alpha\alpha)T=1+2\alpha\cos u+2\alpha\alpha\cos 2u+2\alpha^3\cos 3u+\text{u.s.w.},\tag*{[3]}\end{equation}
\[\begin{aligned}
(1-\alpha\alpha)^3TT={}&1+\alpha\alpha+4\alpha\cos u+(6\alpha\alpha-2\alpha^4)\cos 2u\\
 &+(8\alpha^3-4\alpha^5)\cos 3u+(10\alpha^4-6\alpha^6)\cos 4u+\text{etc.}
\end{aligned}\]
\noindent [Für] \(\alpha=\operatorname{tang}\theta\) [ist]
\begin{equation}T=\frac{1}{1+\alpha\alpha-2\alpha\cos u}=\frac{\cos\theta^2}{1-\sin 2\theta\,{.}\,\cos u},\tag*{[4]}\end{equation}
\begin{equation}(1-\alpha\alpha)T=\frac{\cos 2\theta}{1-\sin 2\theta\,{.}\,\cos u}.\tag*{[5]}\end{equation}
\begin{center}\rule{.24\linewidth}{.4pt}\end{center}
\noindent\hspace*{\parindent}[Es bedeute\textsuperscript{*}) \(v\) die wahre, \(E\) die exzentrische, \(M\) die mittlere Anomalie und \(\varepsilon=\sin\varphi\) die Exzentrizität, so ist]
\par\vfill\noindent\rule{.25\linewidth}{.4pt}\par\smallskip
{\footnotesize\noindent\hspace*{\parindent}[\textsuperscript{*}) Die Bezeichnungen entsprechen denen der \textit{Theoria motus}, {\scriptsize 1809}, siehe Werke VII ({\scriptsize 1906}), S. {\scriptsize 17, 19}.]\par}
\clearpage\SourceLeaf{428}\ArithmeticHead{421}{R}{KONVERGENZ DER ENTWICKLUNG DER MITTELPUNKTSGLEICHUNG.}
\begin{equation}\frac{dv}{dM}=\frac{\surd(1-\varepsilon\varepsilon)}{(1-\varepsilon\cos E)^2}=\frac{\cos\varphi}{(1-\sin\varphi\,{.}\,\cos E)^2}.\tag*{[6]}\end{equation}
\noindent\hspace*{\parindent}[Ferner ist]
\begin{equation}1-\varepsilon\cos E=0\tag*{[7]}\end{equation}
\noindent für
\begin{equation}\left\{\begin{aligned}
\cos E&=\frac1\varepsilon\\
\sin E&=i\frac{\surd(1-\varepsilon\varepsilon)}{\varepsilon}\\
E&=i\log\frac{1+\surd(1-\varepsilon\varepsilon)}{\varepsilon}
\end{aligned}\right.\tag*{[8]}\end{equation}
\begin{equation}M[=E-\varepsilon\sin E]=i\left\{\log\frac{1+\surd(1-\varepsilon\varepsilon)}{\varepsilon}-\surd(1-\varepsilon\varepsilon)\right\}\tag*{[9]}\end{equation}
\begin{equation}\left\{\begin{aligned}
\cos M+i\sin M&=\left\{\frac{1-\surd(1-\varepsilon\varepsilon)}{\varepsilon}\right\}e^{+\surd(1-\varepsilon\varepsilon)}\\
 &=\operatorname{tg}\tfrac12\varphi\,{.}\,e^{+\cos\varphi}.
\end{aligned}\right.\tag*{[10]}\end{equation}
\begin{center}\textls[120]{Beispiel:}\end{center}
\begingroup\small\setlength{\arraycolsep}{8pt}\renewcommand{\arraystretch}{1.35}
\[\begin{array}{lrr}
\hfill\varepsilon=0{,}254236&0{,}4685742&2{,}3666507\\
\hfill\varphi=14^0\,43'\,41''{,}950&\underline{5{,}0172181}&\underline{2{,}3428710}\\
\hfill\cos\varphi\;{.}\,{.}\,{.}\,{.}\,{.}\,{.}\,{.}\,{.}\,{.}\;9{,}9854903&4{,}2115277&1{,}8118966\\
\hfill\kappa[=\log_{10}e]\;{.}\,{.}\,{.}\,{.}\;9{,}6377843&\underline{0{,}9371484}&\underline{2{,}8114452}\\
\hfill\operatorname{tg}\tfrac12\varphi\;{.}\,{.}\,{.}\,{.}\,{.}\,{.}\,{.}\,{.}\;9{,}1114013&3{,}5474520&1{,}2707953\\
\hfill+0{,}4200245&\underline{1{,}4057226}&\underline{3{,}2800194}\\
\raisebox{-.7\baselineskip}{\([\log_{10}\operatorname{tg}\tfrac12\varphi\,{.}\,e^{\cos\varphi}=]\;9{,}5314258\)}&2{,}9412684&0{,}7396514\\
&1{,}8742968&3{,}7485936
\end{array}\]\endgroup
\bigskip\noindent{\small [Zweiter Zettel]}\par
\begin{center}\bfseries Knotenpunkte bei Keplers Problem.\end{center}
\begin{center}[2.]\end{center}
\noindent\hspace*{\parindent}Die Excentricität \(=f\); sonst die gewöhnlichen Bezeichnungen [wie im art. 1.].
\clearpage\SourceLeaf{429}\ArithmeticHead{422}{L}{ANALYSIS. NACHLASS.}
\begin{equation}\left\{\begin{aligned}
\cos E&=\frac1{\sin\varphi}\\
\sin E&=i\operatorname{cotang}\varphi\\
e^{iE}&=\operatorname{tang}\tfrac12\varphi\\
E&=i\log\operatorname{cotg}\tfrac12\varphi
\end{aligned}\right.\tag*{[11]}\end{equation}
\begin{equation}\left\{\begin{aligned}
M&=i(\log\operatorname{cotg}\tfrac12\varphi-\cos\varphi)\\
e^{iM}&=\operatorname{tang}\tfrac12\varphi\,{.}\,e^{\cos\varphi}=\operatorname{tang}\tfrac12\theta\\
\cos M&=\frac1{\sin\theta}=\frac1g\\
\sin M&=i\operatorname{cotg}\theta.
\end{aligned}\right.\tag*{[12]}\end{equation}
\noindent\hspace*{\parindent}Geht hier
\[E\ \text{in}\ E+\varepsilon\]
\noindent und damit
\[M\ \text{in}\ M+\mu\]
\noindent über, so ist
\begin{equation}\mu=\frac12 f\sin E\,{.}\,\varepsilon\varepsilon+\frac16 f\cos E\,{.}\,\varepsilon^3-\frac1{24}f\sin E\,{.}\,\varepsilon^4\ \text{u.s.w.}\tag*{[13]}\end{equation}
\noindent und es werden die Werthe von \(1-f\cos E\)
\begin{equation}=f\sin E\,{.}\,\varepsilon+\frac12 f\cos E\,{.}\,\varepsilon\varepsilon-\frac16 f\sin E\,{.}\,\varepsilon^3\ \text{u.s.w.}\tag*{[14]}\end{equation}
\noindent [und von] \(1-g\cos M\)
\begin{equation}=g\sin M\,{.}\,\mu+\frac12 g\cos M\,{.}\,\mu\mu-\frac16 g\sin M\,{.}\,\mu^3\ \text{u.s.w.}\tag*{[15]}\end{equation}
\noindent\hspace*{\parindent}Man hat also
\begin{equation}\frac{\cos\varphi}{(1-f\cos E)^2}=\frac{g\sin M\,{.}\,\cos\varphi}{2f\sin E\,{.}\,(1-g\cos M)}\Omega,\tag*{[16]}\end{equation}
\noindent wo
\begin{equation}\Omega=\frac{1+\tfrac13\operatorname{cotg}E\,{.}\,\varepsilon}{1+\operatorname{cotg}E\,{.}\,\varepsilon}\tag*{[17]}\end{equation}
\begin{equation}=1-\frac23\operatorname{cotg}E\,{.}\,\varepsilon,\tag*{[18]}\end{equation}
\noindent oder
\begin{equation}\frac{\cos\varphi}{(1-f\cos E)^2}=\frac{\cos\theta}{2(1-g\cos M)}\left(1-\frac23\operatorname{cotg}E\,{.}\,\varepsilon\right).\tag*{[19]}\end{equation}
\noindent Oder, da
\clearpage\SourceLeaf{430}\ArithmeticHead{423}{R}{KONVERGENZ DER ENTWICKLUNG DER MITTELPUNKTSGLEICHUNG.}
\begin{equation}\varepsilon\varepsilon=\frac{2\mu}{f\sin E}=\frac{2(1-g\cos M)}{fg\sin E\,{.}\,\sin M}=\frac{2(1-g\cos M)}{-\cos\varphi\,{.}\,\cos\theta},\tag*{[20]}\end{equation}
\noindent also
\begin{equation}\varepsilon=-i\vsroot{\frac{2(1-g\cos M)}{\cos\varphi\,{.}\,\cos\theta}}\tag*{[21]}\end{equation}
\noindent [ist, so folgt]
\begin{equation}1-\frac23\operatorname{cotg}E\,{.}\,\varepsilon=1+\frac{\surd8}{3\cos\varphi^{\frac32}}\vsroot{\frac{1-g\cos M}{\cos\theta}},\tag*{[22]}\end{equation}
\begin{equation}\frac{\cos\varphi}{(1-f\cos E)^2}=\frac{\cos\theta}{2(1-g\cos M)}+\frac{\surd2}{3\cos\varphi^{\frac32}}\vsroot{\frac{\cos\theta}{1-g\cos M}}.\tag*{[23]}\end{equation}
\begin{center}[3.]\end{center}
\noindent\hspace*{\parindent}Es ist, wenn man \(\operatorname{tang}\tfrac12\theta=\gamma\) setzt [vergl. die Gleichungen [1] bis [5] des art. 1.],
\begin{equation}\left\{\begin{aligned}
\frac{\cos\theta}{1-g\cos M}&=\frac{1-\gamma\gamma}{(1-\gamma e^{iM})(1-\gamma e^{-iM})}\\
 &=1+2\gamma\cos M+2\gamma\gamma\cos 2M+\text{etc.},
\end{aligned}\right.\tag*{[24]}\end{equation}
\begingroup\small
\begin{equation}\begin{gathered}
\vsroot{\frac{\cos\theta}{1-g\cos M}}=\surd(1-\gamma\gamma)\\
\left\{\begin{aligned}
&1+\tfrac12\,{.}\,\tfrac12\gamma\gamma+\tfrac{1\,{.}\,1\,{.}\,3\,{.}\,3}{2\,{.}\,2\,{.}\,4\,{.}\,4}\gamma^4+{.}\,{.}\,{.}\\
&+2\,{.}\,\tfrac12\gamma\left(1+\tfrac12\,{.}\,\tfrac34\gamma\gamma+\tfrac{1\,{.}\,3\,{.}\,3\,{.}\,5}{2\,{.}\,4\,{.}\,4\,{.}\,6}\gamma^4\,{.}\,{.}\,{.}\right)\cos M\\
&+2\,{.}\,\tfrac12\,{.}\,\tfrac34\gamma\gamma\left(1+\tfrac12\,{.}\,\tfrac56\gamma\gamma+{.}\,{.}\,{.}\right)\cos 2M\\
&{.}\quad{.}\quad{.}\quad{.}\quad{.}\quad{.}\quad{.}\quad{.}\quad{.}\quad{.}\quad{.}\quad{.}\quad{.}
\end{aligned}\right.
\end{gathered}\tag*{[25]}\end{equation}
\endgroup
\noindent Also der Coeffic[ient] von \(\cos nM\)
\begin{equation}\begin{aligned}
={}&2\surd(1-\gamma\gamma)\,\frac12\,{.}\,\frac34\,{.}\,\frac56\,{.}\,{.}\,{.}\,\frac{2n-1}{2n}\gamma^n\\
 &\qquad\left(1+\frac12\,{.}\,\frac{2n+1}{2n+2}\gamma\gamma+\text{etc.}\right)
\end{aligned}\tag*{[26]}\end{equation}
\noindent oder proxime
\begin{equation}=2\gamma^n\frac{\Pi2n}{2^{2n}\Pi n^2}=\frac{2\gamma^n}{\surd n\pi}.\tag*{[27]}\end{equation}
\noindent Es ist also in der Entwickl[ung] von \(\frac{dv}{dM}\) [siehe die Gleichung [6] des art. 1.] der Coefficient von \(\cos nM\)
\begin{equation}=\gamma^n\!\left\{1+\frac{\surd8}{3\cos\varphi^{\frac32}\,{.}\,\surd n\pi}\right\}.\tag*{[28]}\end{equation}
\begin{center}\rule{.24\linewidth}{.4pt}\end{center}
\clearpage\endgroup
